\documentclass[10pt,a4paper]{amsart}
\usepackage[margin=19mm]{geometry}
\usepackage{amsmath,amssymb,mathtools}
\usepackage[scr=rsfs,cal=euler]{mathalfa}
\usepackage{xfrac}
\usepackage[unicode,hidelinks,bookmarks=false]{hyperref}
\newcommand{\ie}{i.e.\ }
\newcommand{\cf}{cf.\ }
\newcommand{\resp}{respectively\ }
\newcommand{\Subsection}{\subsection}
\providecommand{\nobreakdash}{\nobreak\hbox{-}}
\setlength{\emergencystretch}{2em}
\multlinegap0pt
\usepackage{amssymb}
\usepackage{mathtools}
\usepackage{bm}
\usepackage[shortlabels]{enumitem}
\usepackage{xcolor}
\newtheorem{proposition}{Proposition}
\newcommand{\dd}{\,d}
\title{\textbf{Critical Compression in Axisymmetric Swirl}}
\author{Rishad Shahmurov}
\date{Source: arXiv:2604.21213v3 (CC BY 4.0)}
\begin{document}
\maketitle

\noindent\emph{Source and licence.} \textbf{Critical Compression in Axisymmetric Swirl}. Version arXiv:2604.21213v3 (CC BY 4.0), distributed by arXiv under the Creative Commons Attribution 4.0 International licence, http://creativecommons.org/licenses/by/4.0/, as recorded in the arXiv licence metadata for this exact version and shown on the abstract page https://arxiv.org/abs/2604.21213v3. Full licence notice: \texttt{licenses/arxiv-2604.21213-CC-BY-4.0.txt}.\par\smallskip

\noindent\textit{Editorial scope.} Counted unit: the article's proposition prop:terminal-no-go (source0.tex lines 862--877). The article's complete original proof of it is delivered with it (source0.tex lines 879--973, one \texttt{proof} environment of 95 lines). No mathematics is written, repaired or re-derived: the statement and the proof are the article's own lines, in the article's own notation, and no retained line is substituted or re-flowed.  No further source-local context is needed: every cross-reference of every retained line resolves inside the two delivered spans.  Nothing was omitted from any retained span, so the package carries no omission record.  No retained line contains a citation, so the wrapper carries no bibliography and no bibliography entry.  No retained line contains a figure environment, an image-inclusion command or drawing code, so no image is delivered and the package declares no asset dependency.  Editorial formatting only, in addition: an AMS \texttt{amsart} preamble that restates the article's own declarations and macros used by the retained lines, namely the environments \texttt{proposition} on separate counters, together with the standard packages those lines need.  The retained environments print in this document as Proposition 1 in order of appearance, whereas the article prints the same items with the numbering of its own full document.  The complete native source of the article as distributed in the arXiv e-print for this exact version is bundled unchanged as \texttt{sources/199006/source0.tex}.
\begin{proposition}\label{prop:terminal-no-go}
A displacement condition imposed only on windows ending at one final time cannot replace the sliding hypothesis.  More precisely, after normalizing $\nu=1$ and $T=1$, for every $K>0$ there exist a nonnegative $g\in C^\infty([0,1))$ and a bounded nonnegative solution of
\[
 q_t=q_{rr}+g(t)q_r,\qquad q(0,t)=0,
\]
with smooth bounded nondecreasing initial data, such that
\begin{equation}\label{eq:terminal-only-wall}
 \int_{1-h}^{1}g(t)\,\dd t
 \le K\sqrt{h\log\log(e^e/h)}
\end{equation}
for every sufficiently small $h$, but there are $t_n\uparrow1$ and $r_n\downarrow0$ for which
\begin{equation}\label{eq:terminal-no-modulus}
 q(r_n,t_n)|\log r_n|^{3/2}\longrightarrow\infty.
\end{equation}
In particular, the terminal-only bound \eqref{eq:terminal-only-wall} does not imply a uniform Wei-type logarithmic modulus on preterminal times.
\end{proposition}

\begin{proof}
Set
\[
 W(h)=\sqrt{h\log\log(e^e/h)}.
\]
Choose $n_0$ large and
\[
 \tau_n=e^{-2^n},\qquad t_n=1-\tau_n,\qquad
 \delta_n=\frac K8 W(\tau_n),\qquad n\ge n_0.
\]
Since $\tau_{n+1}=\tau_n^2$, after increasing $n_0$ we may assume
\begin{equation}\label{eq:deltatail}
 \delta_{n+1}\le\frac14\delta_n,
 \qquad
 \sum_{j\ge n}\delta_j\le\frac43\delta_n.
\end{equation}
Let
\[
 \varepsilon_n=\exp(-\delta_n^{-4}).
\]
Then $\varepsilon_n$ is much smaller than both $\tau_{n+1}$ and the separation $t_{n+1}-t_n$ for large $n$.  Fix $\eta\in C_c^\infty(0,1)$, $\eta\ge0$, $\int_0^1\eta=1$, and put
\[
 g_n(t)=\frac{\delta_n}{\varepsilon_n}
 \eta\!\left(\frac{t-(t_n-\varepsilon_n)}{\varepsilon_n}\right),
 \qquad
 g=\sum_{n\ge n_0}g_n.
\]
The supports are disjoint and locally finite in $[0,1)$, so $g$ is smooth there and each pulse has mass $\int g_n=\delta_n$.

The function $W$ is increasing for all sufficiently small arguments.  Given small $h$, choose $n$ so that $\tau_n\le h<\tau_{n-1}$.  The terminal interval $(1-h,1)$ meets only a part of the $n$th pulse and the full later pulses.  Hence, by \eqref{eq:deltatail},
\[
 \int_{1-h}^1g
 \le\sum_{j\ge n}\delta_j
 \le\frac43\frac K8W(\tau_n)
 \le K W(h),
\]
which proves \eqref{eq:terminal-only-wall}.

It remains to show that each short pulse creates a much thinner preterminal layer than the terminal accounting sees.  Choose smooth bounded increasing initial data
\[
 q_0(r)=\frac{2}{\sqrt\pi}\int_0^r e^{-y^2}\,\dd y.
\]
For each $h>0$, the difference $q(r+h,t)-q(r,t)$ solves the same linear equation on $r>0$ with nonnegative initial data and nonnegative boundary value at $r=0$; the maximum principle therefore gives $q(r+h,t)\ge q(r,t)$.  Hence $q_r\ge0$.  Therefore
\[
 q_t-q_{rr}=gq_r\ge0,
\]
and comparison with the drift-free Dirichlet heat solution $q^{(0)}$ gives $q\ge q^{(0)}$.  By the explicit Dirichlet heat kernel, $\partial_rq^{(0)}(0,t)>0$ for every $t>0$; continuity on the compact interval $[1/2,1]$ yields constants $c_0,r_0>0$ such that
\begin{equation}\label{eq:baseline-linear}
 q(r,t)\ge q^{(0)}(r,t)\ge c_0r,
 \qquad 0<r<r_0,\quad \frac12<t<1.
\end{equation}

Consider the $n$th pulse and set $s_n=t_n-\varepsilon_n$ and
\[
 r_n=4\sqrt{\varepsilon_n}.
\]
During $[s_n,t_n]$ the killed diffusion representation for the time-dependent generator $\partial_{rr}+g(t)\partial_r$ is
\[
 q(r_n,t_n)=
 \mathbb E\left[q(X_{\varepsilon_n},s_n)
 \mathbf 1_{\{\tau_0>\varepsilon_n\}}\right],
\]
where
\[
 X_s=r_n+\int_0^s g(t_n-\xi)\,\dd\xi+\sqrt2 B_s
\]
and $\tau_0$ is the first hitting time of $0$.  The coefficient is read backward from the observation time $t_n$ because this is the Feynman--Kac representation of the forward initial-value problem from $s_n$ to $t_n$.  The deterministic drift is nevertheless nonnegative and has total increment
$\int_0^{\varepsilon_n}g(t_n-\xi)\,\dd\xi=\delta_n$.  Since $\delta_n/\sqrt{\varepsilon_n}\to\infty$, consider the event
\[
 \inf_{0\le s\le\varepsilon_n}\sqrt2B_s\ge-2\sqrt{\varepsilon_n},
 \qquad
 |\sqrt2B_{\varepsilon_n}|\le\frac14\delta_n.
\]
By Brownian scaling and the reflection principle, the first condition has a fixed positive probability, while the probability of the second tends to one.  Their intersection therefore has probability at least a constant $p_0>0$ for all large $n$.  On this event the path remains positive because $r_n=4\sqrt{\varepsilon_n}$ and the deterministic drift is nonnegative.  Moreover, since $4\sqrt{\varepsilon_n}\le\delta_n/4$ for large $n$ and the total deterministic increment equals $\delta_n$,
\[
 \frac12\delta_n\le X_{\varepsilon_n}\le2\delta_n.
\]
For large $n$ this interval lies in $(0,r_0)$.  Applying \eqref{eq:baseline-linear} at time $s_n$ therefore gives
\begin{equation}\label{eq:pulse-lower}
 q(r_n,t_n)\ge c_1\delta_n
\end{equation}
with $c_1>0$ independent of $n$.

Finally,
\[
 |\log r_n|=\frac12\delta_n^{-4}+O(1),
\]
so \eqref{eq:pulse-lower} implies
\[
 q(r_n,t_n)|\log r_n|^{3/2}
 \ge c\delta_n\,\delta_n^{-6}
 =c\delta_n^{-5}\longrightarrow\infty.
\]
This proves \eqref{eq:terminal-no-modulus}.  The normalization $\nu=1$, $T=1$ is removed by the usual parabolic scaling.
\end{proof}

\end{document}
